\section*{Hoofdstuk \ref{ch:b3:bioinformatics} --- Bio-informatica en sequentieanalyse}

\begin{solution}{exo:b3:bioinformatics:1}
Globaal: beide sequenties van begin tot eind uitgelijnd, elk residu in
een kolom --- voor twee eiwitten die over hun hele lengte homoloog
worden geacht, zoals \omterm{def:b3:genomics:comparative}{orthologen} van een huishoudenzym. Lokaal: het
hoogst scorende paar deelreeksen, de rest genegeerd --- om een gedeeld
domein te vinden (een SH2-domein in twee verder ongerelateerde
signaaleiwitten), of een gen in een lange genoomsequentie.
\end{solution}

\begin{solution}{exo:b3:bioinformatics:2}
Randen $0,-1,-2,-3$ beide kanten op. Rij A: $1, 0, -1$. Rij G: $0, 0, -1$.
Rij C: $-1, -1, 1$. Optimum $F(3,3) = 1$: AGC boven AAC zonder hiaten
(overeenkomst, verschil, overeenkomst: $1 - 1 + 1 = 1$).
\end{solution}

\begin{solution}{exo:b3:bioinformatics:3}
$s(a,a) = \lambda^{-1}\log\bigl(q_{aa}/p_{a}^{2}\bigr)$. Tryptofaan is
zeldzaam ($p_{W} \approx 0.013$), dus de kans dat twee tryptofanen
toevallig worden uitgelijnd, $p_{W}^{2}$, is uiterst klein, en een
geconserveerd tryptofaanpaar is een veel sterker teken van homologie
dan een geconserveerd leucinepaar ($p_{L}
\approx 0.1$); de log-oddsverhouding is navenant groter.
\end{solution}

\begin{solution}{exo:b3:bioinformatics:4}
De E-waarde is het aantal \omterm{def:b3:bioinformatics:alignment}{uitlijningen} met een score van minstens die
hoogte dat door toeval te verwachten is bij een zoekactie van deze
zoekvraag tegen een databank van deze omvang. $E = 3$ betekent dat drie
zulke scores door toeval worden verwacht: de treffer is geen bewijs van
homologie (zij kan het nog steeds zijn, maar de zoekactie kan het niet
zeggen).
\end{solution}

\begin{solution}{exo:b3:bioinformatics:5}
$mn = 400\times 2\times 10^{11} = 8\times 10^{13} = 2^{46.2}$. $E(45) =
2^{1.2} \approx 2.3$; $E(55) = 2^{-8.8} \approx 2\times 10^{-3}$;
$E(65) = 2^{-18.8} \approx 2\times 10^{-6}$. Voor $E = 10^{-3}$ is $S' =
46.2 + 10.0 = 56$ bits nodig. Een zoekvraag van \num{40} residuen heeft
een tien keer kleinere $mn$, $2^{42.9}$: $53$ bits volstaan --- maar een
korte zoekvraag haalt zelfs dat zelden.
\end{solution}

\begin{solution}{exo:b3:bioinformatics:6}
\omterm{def:b3:bioinformatics:motif}{Informatie-inhouden}: $2$, $1$, $0$, en $2 - H$ met $H = -(0.7\log_{2}
0.7 + 3\times 0.1\log_{2} 0.1) = 0.36 + 1.00 = 1.36$, dus $0.64$. Totaal
$R = 3.64$ bits. Toevallige treffers: $9.2\times 10^{6}$ posities op
twee strengen $\times 2^{-3.64} \approx 7\times 10^{5}$ --- het \omterm{def:b3:bioinformatics:motif}{motief}
is alleen vrijwel nutteloos.
\end{solution}

\begin{solution}{exo:b3:bioinformatics:7}
Een insertie van verscheidene residuen is één mutatiegebeurtenis, dus
zouden haar kosten niet lineair met haar lengte moeten groeien: een
openingskost plus een kleine verlengingskost bootst dat na. Een te hoge
openingsstraf dwingt verschillen af waar een hiaat hoort en lijnt alles
na een echte insertie verkeerd uit; een te lage straf strooit overal
hiaten, laat residuen toevallig overeenkomen en blaast de identiteit op.
\end{solution}

\begin{solution}{exo:b3:bioinformatics:8}
Niet noodzakelijk: de \omterm{def:b3:bioinformatics:alignment}{uitlijning} beslaat een segment van \num{130}
residuen, en dat is de handtekening van een gedeeld domein en niet van
een \omterm{def:b3:genomics:comparative}{ortholoog} die over zijn hele lengte is uitgelijnd. Toets: doorzoek
met het vliegeneiwit het menselijk eiwitbestand terug (is de zoekvraag
zijn beste treffer, over de hele lengte?), bepaal het domein met een
\omterm{def:b3:bioinformatics:msa}{profiel-HMM}, en bouw een genboom van de familie in verscheidene soorten.
\end{solution}

\begin{solution}{exo:b3:bioinformatics:9}
Een \omterm{def:b3:bioinformatics:msa}{profiel} legt kolom voor kolom vast wat de familie duldt: een
positie die altijd hydrofoob is maar nooit hetzelfde residu, een
onveranderlijk katalytisch residu, een positie die bij de helft van de
familie altijd een hiaat is. Een paarsgewijze \omterm{def:b3:bioinformatics:alignment}{uitlijning} scoort elk
residu tegen één ander residu en kan niet weten dat een valine op
positie 40 ``even goed is als'' het isoleucine dat de zoekvraag daar
heeft. Het \omterm{def:b3:bioinformatics:msa}{profiel} weegt bovendien de geconserveerde kolommen, zodat
zwakke gelijkenis die zich juist daar samenbalt significant wordt.
\end{solution}

\begin{solution}{exo:b3:bioinformatics:10}
Bij een positieve \omterm{prop:b3:bioinformatics:logodds}{verwachte score} $\mu > 0$ per kolom is de cumulatieve
score langs de diagonaal van twee willekeurige sequenties een
toevalswandeling met positieve drift: na $n$ kolommen is zij ongeveer
$\mu n$, zodat de beste \omterm{def:b3:bioinformatics:alignment}{lokale uitlijning} in wezen het geheel is en
haar score als $\mu n$ groeit in plaats van als $\log n$. De theorie
van Karlin--Altschul, die een negatieve drift vergt zodat hoge scores
zeldzame uitstapjes zijn, geldt dan niet en er bestaat geen $\lambda$.
Voor DNA bij \qty{60}{\%} GC is de kans op een overeenkomst
$2(0.3^{2}) + 2(0.2^{2}) = 0.26$, zodat de \omterm{prop:b3:bioinformatics:logodds}{verwachte score}
$0.26 - 0.74 = -0.48$ is: nog steeds negatief, en de statistiek houdt
stand; maar een schema met overeenkomst $+1$ en verschil $-0.3$ zou
verwachting $+0.04$ hebben en zou het hele genoom als één \omterm{def:b3:bioinformatics:alignment}{uitlijning}
melden.
\end{solution}

\begin{solution}{exo:b3:bioinformatics:11}
\omterm{def:b3:bioinformatics:blast}{Seeds} en ketenen: zoek met een hashtabel exacte of bijna exacte
treffers van $k$-meren tussen de twee sequenties, houd die welke op
samenhangende diagonalen liggen, keten ze aaneen, en voer \omterm{thm:b3:bioinformatics:nw}{dynamisch
programmeren} alleen uit in de gaten tussen geketende \omterm{def:b3:bioinformatics:blast}{seeds}; dat geeft
\omterm{def:b3:bioinformatics:alignment}{uitlijningen} op in gebieden zonder \omterm{def:b3:bioinformatics:blast}{seed} (sterk uiteenlopende stukken).
Banden: zijn de twee sequenties bekend als vrijwel collineair, bereken
dan alleen de cellen binnen een band van breedte $w$ rond de diagonaal,
tegen kosten $wn$ in plaats van $mn$; dat geeft elke \omterm{def:b3:bioinformatics:alignment}{uitlijning} op met
een insertie die groter is dan de band.
\end{solution}

\begin{solution}{exo:b3:bioinformatics:12}
Coderende sequentie heeft een periode van drie: de drie codonposities
hebben verschillende basensamenstellingen (de derde is het meest
scheef), en het codongebruik is per soort ongelijk. Een model met drie
coderende toestanden achter elkaar, die elk basen uitzenden met de
samenstelling van die codonpositie, geeft coderend DNA over een venster
van enkele tientallen codons een hogere kans dan de niet-coderende
toestand doet, zelfs zonder stopcodons. In een menselijk genoom zijn de
exons kort (\qty{150}{bp}) en gescheiden door introns van kilobasen,
zodat het coderende signaal kort en onderbroken is; het model moet
bovendien splitsplaatsen herkennen, die zwakke signalen zijn, en de
enorme hoeveelheid niet-coderende sequentie levert vele valse coderende
segmenten op.
\end{solution}

\begin{solution}{pb:b3:bioinformatics:1}
\textbf{1.} Rijen K, Q, T; kolommen K, A, Q, T; randen $0,-1,-2,-3,-4$
en $0,-1,-2,-3$. Rij K: $1, 0, -1, -2$; rij Q: $0, 0, 1, 0$; rij T:
$-1, -1, 0, 2$. Optimum $2$: \texttt{K-QT} boven \texttt{KAQT}.
\textbf{2.} Beste lokale score $3$: \texttt{CAT} tegen \texttt{CAT}
(residuen 4--6 van GATCAT met 2--4 van ACAT); \texttt{ATCAT} tegen
\texttt{A-CAT} scoort eveneens $4 - 1 = 3$.
\textbf{3.} $300\times 450 = 1.35\times 10^{5}$ bewerkingen; tegen de
databank $300\times 1.2\times 10^{11} = 3.6\times 10^{13}$.
\textbf{4.} $3.6\times 10^{4}$ s, tien uur per zoekvraag; de \omterm{def:b3:bioinformatics:blast}{seeds} van
\omterm{def:b3:bioinformatics:blast}{BLAST} slaan vrijwel de hele tabel over en antwoorden in seconden.
\textbf{5.} $q_{ab} = p_{a}p_{b}e^{\lambda s}$. Alanine: $e^{1.39} = 4.0$,
$q_{AA} = 0.074^{2}\times 4.0 = 0.022$, verhouding $4$. Tryptofaan:
$e^{3.82} = 45$, $q_{WW} = 0.013^{2}\times 45 = 0.0077$, verhouding
$45$. Een uitgelijnd tryptofaanpaar komt $45$ keer vaker voor in
homologen dan door toeval, een alaninepaar maar vier keer;
alanineparen zijn in absolute zin niettemin talrijker omdat alanine
veel voorkomt.
\textbf{6.} \qty{24}{\%} ligt in de schemerzone, waar toevallige
\omterm{def:b3:bioinformatics:alignment}{uitlijningen} \qtyrange{15}{20}{\%} halen; de E-waarde van de \omterm{def:b3:bioinformatics:alignment}{uitlijning},
geconserveerde motieven op de juiste posities, een treffer van een
\omterm{def:b3:bioinformatics:msa}{profiel-HMM} op een bekende familie, of een gedeelde vouwing zouden het
beslechten.
\textbf{7.} $mn = 300\times 1.2\times 10^{11} = 3.6\times 10^{13}$;
$\log_{2}(mn) = 45.0$.
\textbf{8.} $E(92) = 2^{45 - 92} = 2^{-47} \approx 7\times 10^{-15}$;
$E(38) = 2^{7} = 128$.
\textbf{9.} $E = 10^{-3}$ bij $S' = 45 + 10 = 55$ bits; $E = 1$ bij $45$
bits.
\textbf{10.} Tien keer $mn$: $E \approx 7\times 10^{-14}$, nog altijd
overweldigend.
\textbf{11.} Met $E = 128$ is de tiende treffer wat het toeval
oplevert; \qty{40}{\%} identiteit over $60$ residuen is geen bewijs. Een
profielzoekactie van de residuen 200--260 tegen de domeindatabank zou
kunnen laten zien of dat segment een bekend domein is, met een
statistiek die de paarsgewijze vergelijking niet heeft.
\textbf{12.} Wederzijds beste treffers over de volle lengte zijn
verenigbaar met een-op-een-orthologie; zij bewijzen haar niet --- een
duplicatie in één lijn na de splitsing geeft twee \omterm{def:b3:genomics:comparative}{co-orthologen}, en het
verlies van de echte \omterm{def:b3:genomics:comparative}{ortholoog} kan een \omterm{def:b3:genomics:comparative}{paraloog} als beste treffer
achterlaten. Een genboom met verscheidene soorten is de toets.
\textbf{13.} $R = 2 + 2 + 1.6 + 2 + 0.8 + 1.2 + 0.4 + 0.3 = 10.3$ bits.
\textbf{14.} $3.2\times 10^{9}\times 2^{-10.3} \approx 2.5\times 10^{6}$
toevallige treffers.
\textbf{15.} $\log_{2}(3.2\times 10^{9}/200) = \log_{2}(1.6\times 10^{7})
\approx 24$ bits.
\textbf{16.} Samen voorkomen op een vaste tussenafstand telt de bits op:
$10.3 + 8 =
18.3$, wat $8$ van de $13.7$ ontbrekende levert; ongeveer $5.7$ bits
(een factor $50$ in toevallige treffers) moet van elders komen ---
toegankelijkheid van het \omterm{def:b3:chromatin-epigenetics:nucleosome}{chromatine}, verdere partners.
\textbf{17.} $H = -(0.5\log_{2}0.5 + 0.5\log_{2}0.5) = 1$ bit; $R = 2 -
1 = 1$ bit.
\textbf{18.} Een bacteriële factor moet haar enkele plaatsen vinden in
een genoom van \qty{4.6}{Mb} zonder hulp van het \omterm{def:b3:chromatin-epigenetics:nucleosome}{chromatine}: zij heeft
zo'n $19$ bits nodig, en haar plaatsen zijn lang en geconserveerd. Een
eukaryoot genoom is duizend keer groter en vraagt tien bits meer, en
toch hebben de factoren korte plaatsen; zij bereiken specificiteit door
combinatie en door de beperking tot toegankelijk \omterm{def:b3:chromatin-epigenetics:nucleosome}{chromatine}, wat de
regulatie ook beter evolueerbaar maakt, omdat een korte plaats
gemakkelijk wordt gewonnen of verloren.
\textbf{19.} $3/64 = 0.047$; het verwachte aantal codons vóór een
stopcodon is $64/3 \approx 21$.
\textbf{20.} $p_{A} = p_{T} = 0.31$, $p_{G} = p_{C} = 0.19$: $P(\text{TAA})
= 0.31^{3} = 0.030$, $P(\text{TAG}) = P(\text{TGA}) = 0.31^{2}\times 0.19
= 0.018$; samen $0.066$, verwachte leesraamlengte $15$ codons. AT-rijk
DNA zit vol stopcodons, dus zijn willekeurige open leesramen korter en
springen lange er sterker uit.
\textbf{21.} $(61/64)^{100} = e^{-4.80} = 0.008$; $(61/64)^{300} =
e^{-14.4} = 5.6\times 10^{-7}$.
\textbf{22.} Elk maximaal open leesraam eindigt op een stopcodon, en
$3.2\times
10^{9}$ codonbeginpunten bevatten $3.2\times 10^{9}\times 3/64 = 1.5\times
10^{8}$ stopcodons: ongeveer $1.5\times 10^{8}\times 0.008 = 1.2\times 10^{6}$
willekeurige leesramen van minstens $100$ codons, en $1.5\times 10^{8}\times
5.6\times 10^{-7} \approx 80$ van minstens $300$.
\textbf{23.} Een bacterie van \qty{4.6}{Mb} heeft zo'n $4\times 10^{5}$
stopcodons en dus ongeveer \num{3500} toevallige leesramen van $100$
codons maar vrijwel geen van $300$; haar genen tellen gemiddeld $300$
codons en \qty{88}{\%} van het DNA is coderend, zodat een lang open
leesraam vrijwel altijd een gen is. Bij de worm codeert \qty{1.5}{\%}
van het DNA, tellen de exons gemiddeld $50$ codons --- korter dan de
toevalsdrempel --- en overspoelen een miljoen toevallige leesramen van
$100$ codons hen. Eukaryote genzoekers gebruiken signalen van
splitsplaatsen, codonvoorkeur in een verborgen markovmodel, homologie
met bekende eiwitten en, vooral, gesequencede transcripten.
\textbf{24.} Een \omterm{def:b3:genomics:ngs}{read} uit een gesplitste boodschapper lijnt in twee
stukken op het genoom uit, gescheiden door een intron: de splitsing
markeert beide splitsplaatsen tot op de base; de readdekking bakent de
exons af en gepaarde \omterm{def:b3:genomics:ngs}{reads} verbinden opeenvolgende exons tot één
transcript, wat uitwijst welke van verscheidene kandidaat-splitsplaatsen
wordt gebruikt.
\textbf{25.} $E \approx 7\times 10^{-15}$ voor de beste treffer;
ongeveer $24$ bits om $200$ plaatsen in het genoom aan te wijzen; zo'n
$80$ toevallige leesramen van $300$ codons in het hele genoom.
\end{solution}
